\documentclass[12pt]{article}

\usepackage[a4paper,margin=1in]{geometry}
\usepackage{amsmath,amssymb}
\usepackage{setspace}

\setstretch{1.2}

\title{CSE400 – Fundamentals of Probability in Computing\\
Lecture 19: N-Variate Joint PDF/PMF/CDF}

\date{}

\begin{document}

\maketitle

\section*{1. Joint Distribution of N Random Variables}

\subsection*{1.1 Representation of Multiple Random Variables}

\textbf{Definition}

For random variables $X_1, X_2, \dots, X_N$, the vector form is:

\[
X =
\begin{bmatrix}
X_1 \\
X_2 \\
\vdots \\
X_N
\end{bmatrix}
\]

The joint distribution describes the probability behavior of all variables together.

It can be represented using:

Joint PMF

Joint PDF

Joint CDF

\section*{2. Jointly Gaussian Random Variables}

\subsection*{2.1 Joint Gaussian PDF (Two Variables)}

For two variables $X$ and $Y$:

\[
f_{XY}(x,y)=
\frac{1}{2\pi\sigma_X\sigma_Y\sqrt{1-\rho^2}}
\exp\left(
-\frac{1}{2(1-\rho^2)}
\left[
\left(\frac{x-\mu_X}{\sigma_X}\right)^2
+
\left(\frac{y-\mu_Y}{\sigma_Y}\right)^2
-2\rho
\left(\frac{x-\mu_X}{\sigma_X}\right)
\left(\frac{y-\mu_Y}{\sigma_Y}\right)
\right]
\right)
\]

Where:

$\mu_X,\mu_Y$: means

$\sigma_X,\sigma_Y$: standard deviations

$\rho$: Pearson correlation coefficient

\section*{3. Example 1: Jointly Gaussian Variables}

Given

\[
\mu_X=2,\quad \sigma_X^2=1
\]

\[
\mu_Y=-3,\quad \sigma_Y^2=4
\]

\[
cov(X,Y)=-1
\]

\subsection*{(a) Joint PDF Construction}

Step 1: Correlation Coefficient

\[
\rho_{XY}=\frac{cov(X,Y)}{\sigma_X\sigma_Y}
\]

\[
\sigma_X=1,\quad \sigma_Y=2
\]

\[
\rho_{XY}=-\frac{1}{2}
\]

Step 2: Substitute into General Formula

\[
\mu_X=2,\ \sigma_X=1,\ \mu_Y=-3,\ \sigma_Y=2,\ \rho=-0.5
\]

Final Joint PDF

\[
f_{XY}(x,y)=
\frac{1}{2\pi\sqrt{3}}
\exp\left(
-\frac{2}{3}
\left[
(x-2)^2
+
(x-2)(y+3)
+
\frac{(y+3)^2}{4}
\right]
\right)
\]

\subsection*{(b) Marginal PDFs}

\[
f_X(x)=\int_{-\infty}^{\infty} f_{XY}(x,y) dy
\]

\[
f_Y(y)=\int_{-\infty}^{\infty} f_{XY}(x,y) dx
\]

Key Property

If $X$ and $Y$ are jointly Gaussian:

\[
X\sim N(\mu_X,\sigma_X^2),\quad Y\sim N(\mu_Y,\sigma_Y^2)
\]

Marginal Results

\[
f_X(x)=\frac{1}{\sqrt{2\pi}} \exp\left(-\frac{(x-2)^2}{2}\right)
\]

\[
f_Y(y)=\frac{1}{2\sqrt{2\pi}} \exp\left(-\frac{(y+3)^2}{8}\right)
\]

\subsection*{(c) Probability $Pr(X<0)$}

Step 1: Identify CDF

\[
F_X(x)=Pr(X\le x),\quad Pr(X<0)=F_X(0)
\]

Step 2: Standardization

\[
Z=\frac{X-\mu_X}{\sigma_X}
\]

\[
Pr(X<0)=Pr(Z<-2)
\]

Step 3: Symmetry Property

\[
\phi(-x)=1-\phi(x)=Pr(Z>x)
\]

\[
Pr(Z<-2)=Pr(Z>2)
\]

Step 4: Q-Function

\[
Q(x)=Pr(Z>x)
\]

Final Answer

\[
Pr(X<0)=Q(2)
\]

\section*{4. Application: Large MIMO System}

System Setup

$M$: transmit antennas

$N$: receive antennas

Channel Coefficients

Each path between antenna $i$ and $j$:

$h_{ij}$

Channel Matrix

\[
H =
\begin{bmatrix}
h_{11} & h_{12} & \cdots & h_{1N} \\
h_{21} & h_{22} & \cdots & h_{2N} \\
\vdots & \vdots & \ddots & \vdots \\
h_{M1} & h_{M2} & \cdots & h_{MN}
\end{bmatrix}_{M\times N}
\]

\section*{5. Joint PMF}

For discrete random variables:

\[
P_{X_1,\dots,X_N}(x_1,\dots,x_N)=
P(X_1=x_1,\dots,X_N=x_N)
\]

\section*{6. Joint CDF}

\[
F_{X_1,\dots,X_N}(x_1,\dots,x_N)=
P(X_1\le x_1,\dots,X_N\le x_N)
\]

\section*{7. Joint PDF}

\[
f_{X_1,\dots,X_N}(x_1,\dots,x_N)=
\frac{\partial^N}{\partial x_1\cdots \partial x_N}
F_{X_1,\dots,X_N}(x_1,\dots,x_N)
\]

\section*{8. Independence of Random Variables}

Condition

\[
X=[X_1,X_2,\dots,X_N]^T
\]

If independent:

\[
f_X(x_1,\dots,x_N)=\prod_{i=1}^N f_{X_i}(x_i)
\]

\section*{9. Example 2: Uniformly Distributed 3D Vector}

Given

\[
|X_1|\le1,\ |X_2|\le1,\ |X_3|\le1
\]

\[
f_{X_1,X_2,X_3}(x_1,x_2,x_3)=
\begin{cases}
C, & \text{inside cube} \\
0, & \text{otherwise}
\end{cases}
\]

(a) Find $C$

\[
C(2)(2)(2)=1 \Rightarrow C=\frac{1}{8}
\]

(b) Marginal

\[
f_{X_1,X_2}(x_1,x_2)=\int_{-1}^{1}\frac{1}{8}dx_3=\frac{1}{4}
\]

(c) Marginal

\[
f_{X_1}(x_1)=\int_{-1}^{1}\int_{-1}^{1}\frac{1}{8}dx_2dx_3=\frac{1}{2}
\]

(d) Conditional PDF

\[
f_{X_1|X_2,X_3}=\frac{1/8}{1/4}=\frac{1}{2}
\]

(e) Conditional PDF

\[
f_{X_1,X_2|X_3}=\frac{1/8}{1/2}=\frac{1}{4}
\]

(f) Independence Check

\[
LHS=\frac{1}{8}
\]

\[
RHS=\frac{1}{2}\cdot\frac{1}{2}\cdot\frac{1}{2}=\frac{1}{8}
\]

Conclusion

LHS = RHS $\Rightarrow$ independent

\section*{10. Example 3: Trivariate Random Variables}

Given

\[
f_{XYZ}(x,y,z)=
\begin{cases}
Ke^{-(ax+by+cz)}, & x,y,z>0 \\
0, & \text{otherwise}
\end{cases}
\]

(a) Find $K$

\[
K\left(\frac{1}{a}\right)\left(\frac{1}{b}\right)\left(\frac{1}{c}\right)=1
\]

\[
K=abc
\]

(b) Marginal Joint PDF

\[
f_{XY}(x,y)=\int_0^\infty abce^{-(ax+by+cz)}dz
=abe^{-(ax+by)}
\]

(c) Marginal PDF

\[
f_X(x)=\int_0^\infty abe^{-(ax+by)}dy
=ae^{-ax}
\]

(d) Independence Check

LHS = $abce^{-(ax+by+cz)}$

RHS = $(ae^{-ax})(be^{-by})(ce^{-cz})$

Conclusion

LHS = RHS $\Rightarrow X,Y,Z$ are independent

\end{document}
